% Supplementary material for Section~\ref{sec:mechanism} (Does AWWL change the
% spectrum?). Input under Appendix D (Supplementary results); starts at
% subsection level. Moved floats keep their labels: tab:bands-ddpm,
% fig:compare.

\subsection{The spectrum under DDPM-1000, and paired samples}
\label{app:supp-mechanism}

Table~\ref{tab:bands-ddpm} repeats Table~\ref{tab:bands} for ancestral
DDPM-1000 sampling of the same weights (Section~\ref{sec:mechanism}). Under
this sampler the baseline's deficit is smaller in every band, the static
wavelet arm stays close to the MSE level, as under DDIM-100 ($\corrHighStaticddpm$~dB
$[\corrHighStaticddpmLo, \corrHighStaticddpmHi]$,
$p = \corrHighStaticddpmP$), and AWWL's high-band change is a small shift
towards a larger deficit ($\corrHighAWWLddpm$~dB
$[\corrHighAWWLddpmLo, \corrHighAWWLddpmHi]$, $p = \corrHighAWWLddpmP$).

\IfFileExists{tables/bands_ddpm1000.tex}{% Generated by scripts/make_band_table.py --- do not edit by hand.
% Recomputed from the sampled images (all samples, all reference images).

\begin{table}[tbp]
\centering
\footnotesize
\setlength{\tabcolsep}{3pt}
\caption{Signed deviation of the samples' radial luminance spectrum from the reference set's, in dB, at epoch 199 (DDPM-1000); negative means too little energy. Every sample of each run and all \num{50000} CIFAR-10 training images enter the profiles. Bands are thirds of the integer radial bins: radii 0--6 (low), 7--14 (mid) and 15--22 (high), in cycles per image, where the axial Nyquist radius is 16. A band's value is the unweighted mean of its bins' deviations, so each radius counts once whatever its number of coefficients. Mean $\pm$ sample standard deviation over five seeds. The last two columns pair each seed with its own MSE run in the high band: the difference of deviations (positive means more high-frequency energy than MSE), its 95\% $t$-interval and the two-sided paired $t$-test $p$, uncorrected.}
\label{tab:bands-ddpm}
\begin{tabular}{lrrrrr}
\toprule
Configuration & Low & Mid & High & $\Delta$ high [95\% CI] & $p$ \\
\midrule
MSE & $-0.14 \pm 0.05$ & $-0.42 \pm 0.04$ & $-1.30 \pm 0.04$ & --- & --- \\
AWWL & $-0.13 \pm 0.04$ & $-0.41 \pm 0.04$ & $-1.34 \pm 0.03$ & $-0.04$\;$[-0.06, -0.03]$ & $0.0021$ \\
Static wavelet ($p = 0$) & $-0.14 \pm 0.04$ & $-0.42 \pm 0.04$ & $-1.32 \pm 0.05$ & $-0.02$\;$[-0.05, 0.004]$ & $0.074$ \\
\bottomrule
\end{tabular}
\end{table}
}{%
  \begin{table}[tbp]\centering
  \fbox{\textbf{tables/bands\_ddpm1000.tex not built}}
  \caption{Not built.}\label{tab:bands-ddpm}
  \end{table}}

Figure~\ref{fig:compare} pairs uncurated samples of MSE, the static wavelet
arm and AWWL drawn from the same initial noise. The difference row shows small
differences spread over the whole image and larger ones concentrated at object
boundaries and in small structures. Across the three sample rows, no
systematic difference is apparent on inspection at print size.

\guardedfigure{1.0}{compare.png}{Uncurated samples (indices 0--7) of training
seed 1 at epoch 199, drawn with DDIM-100 ($\eta = 0$). All rows of a column
start from the same initial noise, and DDIM with $\eta = 0$ is deterministic
given that noise, so the images in a column differ only through the training
objective. The last row is $\lvert \mathrm{AWWL} - \mathrm{MSE} \rvert$,
averaged over the RGB channels and amplified eight-fold (darker is larger;
differences of 32 or more grey levels print black).}{fig:compare}
